\input{../tex-inputs/latex-leseansicht-vorspann}

\section[Arthur und Olga Schnitzler an Gustav Schwarzkopf, 1{[}1{]}. 6. 1917]{L04411 Arthur und Olga Schnitzler an Gustav Schwarzkopf, 1{[}1{]}. 6. 1917}
\nopagebreak\mylabel{L04411v}
\rehead{ }\normalsize\beginnumbering\briefempfaengerindex{Schwarzkopf, Gustav@\textsc{Schwarzkopf, Gustav}!zzzSchnitzler, Olga@\emph{von Olga Schnitzler}!1917-06-112@{1{[}1{]}. 6. 1917}|(be}\briefempfaengerindex{Schwarzkopf, Gustav@\textsc{Schwarzkopf, Gustav}!zzzSchnitzler, Arthur@\emph{von Arthur Schnitzler}!1917-06-112@{1{[}1{]}. 6. 1917}|(be}
\toendnotes[C]{\smallbreak\pagebreak[2]}
\correspDesc{Versand  durch Arthur Schnitzler, Olga Schnitzler am 1[1]. 6. 1917 in Bad Gastein
\newline{}Übermittlung  am 11. 6. 1917 in Bad Gastein
\newline{}Erhalt  durch Gustav Schwarzkopf im Zeitraum [12. 6. 1917 – 16. 6. 1917?] in Wien}\toendnotes[C]{\smallbreak}
\Standort{DLA, A:Schnitzler, HS.1985.1.1897.}
\physDesc{Bildpostkarte, 676 Zeichen
\newline{}Handschrift Arthur Schnitzler: Bleistift, lateinische Kurrent
\newline{}Handschrift Olga Schnitzler: Bleistift, lateinische Kurrent
\newline{}Versand: Stempel: »\nobreak{}\oindex{Bad Gastein@\textbf{Bad Gastein}|pwk}Bad
                                       Gaste\textcolor{gray}{in}, 11. VI. 17, 5\nobreak{}«.  }\toendnotes[C]{\smallbreak}\pstart{}{\pb}Herrn Gustav Schwarzkopf\pend{}\pstart{}Wien I\oindex{I., Innere Stadt@\textbf{I., Innere Stadt}|pw}\pend{}\pstart{}Tiefer Graben 17\oindex{Wien@\textbf{Wien}!I., Innere Stadt@\textbf{I., Innere Stadt}!Tiefer Graben 17@\textbf{Tiefer Graben 17}|pw}\pend{}{\bigskip}
\pstart
           \centering{}{\pb}\textcolor{gray}{\textbf{BAD GASTEIN\oindex{Bad Gastein@\textbf{Bad Gastein}|pw}}}\pend
           \vspace{1em}
\pstart
           \noindent{}{\pb}lieber Gustav – wir grüßen Sie
               herzlichst  und theilen mit daſs wir uns bei etwas wechselndem, meist schönem Wetter,
               vortrefflichem Quartier und vorzüglichem (we{\geminationn} auch nicht
               allzu reichlichem) Essen recht wohl befinden u es zu verschmerzen versuchen, daß wir
               der \label{K_L04371-1v}\edtext{ersten Première\eventindex{Burgtheater@\textbf{Burgtheater}!Premiere von Die goldene Eva, 16.6.1917@Premiere von Die goldene Eva, 16.6.1917|pwv} unter dem neuen Herrn\pwindex{Millenkovich, Max von 2.\,3.\,1866 Wien – 5.\,2.\,1945 Baden bei Wien@\textsc{Millenkovich, Max von} (2.\,3.\,1866 Wien – 5.\,2.\,1945 Baden bei Wien)|pwv}}{\lemma{\textnormal{\emph{ersten … Herrn}}}\Cendnote{\textnormal{Schnitzler witzelt
                  darüber, dass der seit April 1917 als \emph{Burgtheaterdirektor}\orgindex{Burgtheater@Burgtheater|pwk} amtierende Max von
                     Millenkovich\pwindex{Millenkovich, Max von 2.\,3.\,1866 Wien – 5.\,2.\,1945 Baden bei Wien@\textsc{Millenkovich, Max von} (2.\,3.\,1866 Wien – 5.\,2.\,1945 Baden bei Wien)|pwk} als erste Premiere\eventindex{Burgtheater@\textbf{Burgtheater}!Premiere von Die goldene Eva, 16.6.1917@Premiere von Die goldene Eva, 16.6.1917|pwk} das
                  gefällige Lustspiel \emph{Die goldene Eva}\pwindex{Schönthan-Pernwald, Franz von 20.\,6.\,1849 Wien – 2.\,12.\,1913 ebd.@\textsc{Schönthan-Pernwald, Franz von} (20.\,6.\,1849 Wien – 2.\,12.\,1913 ebd.)!goldene Eva. Lustspiel in drei Akten@\strich\emph{Die goldene Eva. Lustspiel in drei Akten}|pwk}\pwindex{Koppel-Ellfeld, Franz 7.\,12.\,1838 Eltville – 16.\,1.\,1920 Dresden@\textsc{Koppel-Ellfeld, Franz} (7.\,12.\,1838 Eltville – 16.\,1.\,1920 Dresden)!goldene Eva. Lustspiel in drei Akten@\strich\emph{Die goldene Eva. Lustspiel in drei Akten}|pwk}
                  ansetzte, für den 16. 6. 1917. }}}\label{K_L04371-1} (Junifestspiele
               sollen es wohl sein –?)  nicht beiwohnen können. Lassen Sie bald hören, wies
               Ihnen geht. Noch ziemlich leer ist es hier –  {\pb}man
               ruht sich so recht aus – und geht viel spazieren.\pend
           
\pstart
           Ihr{\\[\baselineskip]}\spacefill\mbox{Arthur}\pend
           \leftskip=0em{}
\pstart
           \noindent{}Viele Grüße Ihnen Beiden\pwindex{Schwarzkopf, Max 12.\,6.\,1857 Wien – 14.\,4.\,1928 ebd.@\textsc{Schwarzkopf, Max} (12.\,6.\,1857 Wien – 14.\,4.\,1928 ebd.)|pwv}.\pend
           \selectlanguage{ngerman}\vspace{1em}
\pstart
           \noindent{}{[}hs. O. Schnitzler:{]} Hab ich nicht Pech? endlich wird sie aufgeführt\eventindex{Burgtheater@\textbf{Burgtheater}!Premiere von Die goldene Eva, 16.6.1917@Premiere von Die goldene Eva, 16.6.1917|pw}, die goldene
                  Eva\pwindex{Schönthan-Pernwald, Franz von 20.\,6.\,1849 Wien – 2.\,12.\,1913 ebd.@\textsc{Schönthan-Pernwald, Franz von} (20.\,6.\,1849 Wien – 2.\,12.\,1913 ebd.)!goldene Eva. Lustspiel in drei Akten@\strich\emph{Die goldene Eva. Lustspiel in drei Akten}|pw}\pwindex{Koppel-Ellfeld, Franz 7.\,12.\,1838 Eltville – 16.\,1.\,1920 Dresden@\textsc{Koppel-Ellfeld, Franz} (7.\,12.\,1838 Eltville – 16.\,1.\,1920 Dresden)!goldene Eva. Lustspiel in drei Akten@\strich\emph{Die goldene Eva. Lustspiel in drei Akten}|pw} und ich kann nicht dabei sein!\pend
           \pstart Herzlichst Ihre \spacefill\mbox{O. S.}\pend{}\selectlanguage{ngerman}\endnumbering\briefempfaengerindex{Schwarzkopf, Gustav@\textsc{Schwarzkopf, Gustav}!zzzSchnitzler, Olga@\emph{von Olga Schnitzler}!1917-06-112@{1{[}1{]}. 6. 1917}|)be}\briefempfaengerindex{Schwarzkopf, Gustav@\textsc{Schwarzkopf, Gustav}!zzzSchnitzler, Arthur@\emph{von Arthur Schnitzler}!1917-06-112@{1{[}1{]}. 6. 1917}|)be}\mylabel{L04411h}
\begin{anhang}
\end{anhang}\newcommand{\dateiname}{L04411}\newcommand{\titel}{Arthur und Olga Schnitzler an Gustav Schwarzkopf, 1[1]. 6. 1917}\newcommand{\editorInnen}{Herausgegeben von Jahnke, SelmaMüller, Martin Anton}\input{../tex-inputs/latex-leseansicht-abspann}
